\section{Project Organization and Repository Structure}
\label{sec:project-organization}

% Section 01.04: Project Organization and Repository Structure

A well-organized repository structure is essential for maintainability, especially in embedded projects where you may support multiple chip variants, boards, and protocol stacks. In this section, we'll examine the repository structure used throughout this book and understand the rationale behind each decision.

\subsection{Top-Level Directory Structure}

Here's the complete structure of our EFR32MG24 project:

% Directory Tree for Section 1.4 - Repository Structure
\begin{figure}[H]
    \centering
    \begin{forest}
        for tree={
            font=\ttfamily\footnotesize,
            grow'=0,
            child anchor=west,
            parent anchor=south,
            anchor=west,
            calign=first,
            edge path={
                \noexpand\path [draw] (!u.south west) +(7.5pt,0) |- node[fill,inner sep=1.25pt] {} (.child anchor);
            },
            before typesetting nodes={
                if n=1
                    {insert before={[,phantom]}}
                    {}
            },
            fit=band,
            before computing xy={l=15pt},
        }
        [EFR32MG24/
            [efr32mg24-pac/
                [Cargo.toml]
                [src/]
                [svd/]
            ]
            [efr32mg24-hal/
                [Cargo.toml]
                [src/]
                [examples/]
            ]
            [efr32mg24-rail-sys/]
            [efr32mg24-bluetooth-sys/]
            [xiao-mg24-bsp/]
            [docs/]
            [book/
                [latex/]
                [TOC.yml]
            ]
            [examples/]
            [Cargo.toml]
            [README.md]
            [CLAUDE.md]
        ]
    \end{forest}
    \caption{EFR32MG24 project repository structure}
    \label{fig:repo-structure}
\end{figure}


\subsection{Crate Organization}

\subsection{Peripheral Access Crate (efr32mg24-pac)}

\begin{minted}{text}
efr32mg24-pac/
├── Cargo.toml          # PAC crate metadata
├── build.rs            # Build script (if needed)
├── svd/
│   └── EFR32MG24B220F1536IM48.svd  # SVD file
├── src/
│   └── lib.rs         # 138,000+ lines of generated code
└── README.md          # PAC-specific documentation
\end{minted}

**Key Points**:
\begin{itemize}
    \item Contains ONLY auto-generated code from svd2rust
    \item Never edited manually after generation
    \item Version locked to specific SVD file version
    \item Re-exported by HAL, users never depend on PAC directly
\end{itemize}

\subsection{Hardware Abstraction Layer (efr32mg24-hal)}

The HAL workspace contains multiple peripheral crates:

\begin{minted}{text}
efr32mg24-hal/
├── Cargo.toml              # Workspace root
├── src/
│   ├── lib.rs             # Main HAL entry point
│   ├── clock/             # Clock management (CMU)
│   │   ├── mod.rs
│   │   ├── types.rs
│   │   ├── traits.rs
│   │   └── README.md
│   ├── gpio/              # GPIO peripheral
│   │   ├── mod.rs
│   │   ├── types.rs
│   │   ├── traits.rs
│   │   └── README.md
│   ├── usart/             # USART communication
│   │   ├── mod.rs
│   │   ├── types.rs
│   │   ├── traits.rs
│   │   └── README.md
│   ├── i2c/               # I2C peripheral
│   │   ├── mod.rs
│   │   ├── types.rs
│   │   ├── traits.rs
│   │   └── README.md
│   ├── spi/               # SPI peripheral
│   │   ├── mod.rs
│   │   ├── types.rs
│   │   ├── traits.rs
│   │   └── README.md
│   └── prelude.rs         # Common imports
├── examples/              # HAL-level examples
│   ├── 01_clock.rs
│   ├── 02_gpio_blink.rs
│   ├── 03_usart_echo.rs
│   └── ...
└── docs/                  # HAL-specific documentation
    ├── README.md
    └── BUILDING.md
\end{minted}

\textbf{File Pattern}: Each peripheral follows the same structure with four files: mod.rs, types.rs, traits.rs, and README.md.

\begin{minted}{text}
peripheral-name/
├── mod.rs       // Main implementation
├── types.rs     // Type definitions
├── traits.rs    // Trait implementations
└── README.md    // Per-peripheral documentation
\end{minted}

\subsection{Radio System Crates (efr32mg24-rail-sys)}

FFI bindings to Silicon Labs RAIL (Radio Abstraction Interface Layer):

\begin{minted}{text}
efr32mg24-rail-sys/
├── Cargo.toml          // Link to vendor libraries
├── build.rs
├── wrapper.h          // bindgen header
└── src/
    └── lib.rs         # Generated FFI bindings
\end{minted}

**Purpose**: Enables multi-protocol radio operation (Bluetooth, Zigbee, Thread, proprietary).

\subsection{Protocol System Crates (bluetooth-sys, etc.)}

FFI bindings to higher-level protocol stacks:

\begin{minted}{text}
efr32mg24-bluetooth-sys/
├── Cargo.toml         # Bluetooth library linking
├── build.rs
├── wrapper.h
└── src/
    ├── lib.rs        # FFI bindings
    └── events/       # Event handling code
\end{minted}

\subsection{Board Support Package (xiao-mg24-bsp)}

Board-specific abstractions:

\begin{minted}{text}
xiao-mg24-bsp/
├── Cargo.toml          # Depends on efr32mg24-hal
├── src/
│   ├── lib.rs         # Board initialization
│   ├── board.rs       # Board struct
│   ├── sensors.rs     # Onboard sensors
│   └── radio.rs       # Radio utilities
└── examples/
    ├── blink_onboard_led.rs
    ├── ble_beacon.rs
    └── zigbee_coordinator.rs
\end{minted}

\subsection{Why This Structure Works}

\subsection{Clear Separation of Concerns}
\begin{itemize}
    \item PAC: Hardware truth (auto-generated)
    \item HAL: Safe abstractions (manually written)
    \item BSP: Board integration (thin wrappers)
    \item Application: User logic (no hardware details)
\end{itemize}

\subsection{Testing Strategy}
Each layer can be tested independently:
\begin{itemize}
    \item PAC: Verify register addresses match datasheet
    \item HAL: Unit tests with mocked registers
    \item BSP: Integration tests on real hardware
    \item Application: Standard Rust testing
\end{itemize}

\subsection{Version Management}
Crates can be versioned independently:
\begin{minted}{text}
efr32mg24-pac = "0.1.0"      # Hardware revision specific
efr32mg24-hal = "0.4.2"      # Feature updates
xiao-mg24-bsp = "0.2.1"      # Board revision specific
\end{minted}

\begin{exercise}[Easy]
    Sketch a project structure for supporting three boards:
    \begin{itemize}
        \item Seeed XIAO MG24 (covered here)
        \item Silicon Labs xG24-DK2601B dev kit
        \item Custom PCB with external sensors
    \end{itemize}
    What crates would be shared? Which would be board-specific?
\end{exercise}

\begin{exercise}[Medium]
    Consider a startup shipping products using this architecture:
    \begin{enumerate}
        \item How would you version the crates when releasing new firmware?
        \item What testing strategy ensures HAL changes don't break BSP?
        \item How would you handle multiple hardware revisions?
    \end{enumerate}
\end{exercise}

\begin{exercise}[Hard]
    Compare this structure to Linux kernel architecture:
    \begin{enumerate}
        \item Map PAC/HAL/BSP to Linux equivalents
        \item How does Linux handle different architectures? (arch/ directory)
        \item How does Linux handle different boards? (Device Tree)
        \item What lessons can we apply to our Rust embedded architecture?
    \end{enumerate}
\end{exercise}
\clearpage
